% !PRINTED 34
\weberpage{34}

If a function is divisible by another, then the degree of the quotient is equal to the difference between the degree of the dividend and the degree of the divisor. The former therefore cannot be smaller than the latter. If the degrees are equal, the quotient is a constant, and hence:

\gesperrt{7. If, of two whole rational functions of the same degree, one is divisible by the other, then they differ from one another only by a constant factor, and the second is also divisible by the first.}

\gesperrt{8. By \S~4, the necessary and sufficient condition that the function $f(x)$ be divisible by the linear function $x-\alpha$ is that $f(\alpha)=0$.}

\begin{center}
{\large \S.\ 6.}

\vspace{0.5em}
{\large\gesperrt{Greatest Common Divisor.}}
\end{center}

It is a problem of fundamental importance to decide whether two whole rational functions have, besides the constants, any further common divisor. The solution is found by the \gesperrt{algorithm of the greatest common divisor} in exactly the same way as the corresponding question concerning divisibility of integers is answered. See the Introduction.

Let
\begin{equation}
 f(x)=A,
 \qquad
 \varphi(x)=A'
\tag{1}
\end{equation}
be two given functions; the degree of $\varphi(x)$ is supposed to be lower, or at least not higher, than that of $f(x)$.

We divide $A$ by $A'$ and denote the remainder, whose degree is lower than the degree of $A'$, by $A''$, thus:
\[
 A=Q'A'+A'';
\]
then we divide $A'$ by $A''$ and denote the remainder by $A'''$, and continue in this way in forming the sequence of functions
\begin{equation}
 A,A',A'',A''',\ldots,
\tag{2}
\end{equation}
whose degrees $n,n',n'',n'''\ldots$ always decrease and therefore, after a finite number of divisions, necessarily go down to zero. Let the last remainder of degree zero, and therefore a constant, be $A^{(\nu)}$. Then we have the chain of equations:
%CONT>>
